\chapter{Arnés de paridad}
\label{ch:parity-harness}

\begin{note}
Stub. El arnés vive en
\filepath{graph\_hunter\_api/tests/parity*.rs}. Dos fases se
entregan hoy (A: forma de DTO; B: escenario); una tercera
cubriendo el runtime HTTP + Tauri queda diferida.
\end{note}

\section{Fase A --- forma de DTO}

Para cada par de DTO request/response, el arnés serializa un
valor armado a mano a JSON y asevera que las claves top-level
del objeto coincidan exactamente con un conjunto esperado. Un
campo que falta o un rename levanta un error claro en test time.

El helper:

\begin{lstlisting}[style=rust]
fn assert_keys<T: Serialize>(value: &T, expected: &[&str]) {
    let v = serde_json::to_value(value).unwrap();
    let actual: BTreeSet<_> = v.as_object().unwrap()
        .keys().map(|s| s.as_str()).collect();
    let want: BTreeSet<_> = expected.iter().copied().collect();
    assert_eq!(actual, want);
}
\end{lstlisting}

\section{Fase B --- escenario}

Los escenarios end-to-end ejercitan secuencias ordenadas de
operaciones contra la fachada directamente (no a través de un
transport). El arnés construye una sesión fixture con un
mini-grafo seed (Alice + Bob autenticándose a web-01, web-01
ejecutando cmd.exe), corre una secuencia de llamadas de API y
asevera estado observable en cada paso.

Fixtures:
\begin{itemize}
    \item \code{fresh\_api\_with\_session()} --- API fresca vacía
        más una sesión única.
    \item \code{fresh\_api\_with\_seeded\_session()} --- el grafo
        seed de 4 entidades / 3 aristas.
    \item \code{tag}, \code{note}, \code{parse} --- helpers de
        conveniencia para que los escenarios se lean al nivel
        de operación.
\end{itemize}

Aserciones de muestra:
\code{run\_hunt\_populates\_cache\_and\_pages\_back} (un run
completo $\to$ round-trip de page-back matchea),
\code{diff\_hunts\_rejects\_inverted\_window} (el invariante del
Capítulo~\ref{ch:hunt-diff}),
\code{export\_iocs\_sentinel\_watchlist\_csv\_contains\_tagged\_ips\_only}
(el formato de export).

\section{Fase C diferida}

Un arnés runtime HTTP + Tauri completo (levantar axum en el
test, cablear un mock de IPC Tauri) está marcado como follow-up.
Fase A + B juntas atrapan los drifts de forma y semántica que
habrían roto la paridad entre transports; Fase C agregaría
round-trips a nivel de wire.

\begin{inthecode}
\filepath{graph\_hunter\_api/tests/parity.rs} --- entry del arnés.\\
\filepath{graph\_hunter\_api/tests/parity/fixture.rs} ---
constructores de fixtures.\\
\filepath{graph\_hunter\_api/tests/parity/dto\_shapes.rs} --- Fase A.\\
\filepath{graph\_hunter\_api/tests/parity/scenarios.rs} --- Fase B.
\end{inthecode}
